\section{Reproducibility and AI-Assisted Workflow}
\paragraph{Source and execution.}
The author WITT implementation is pinned to commit \texttt{2460db5} at \url{https://github.com/KeYang8/WITT}. Evaluation records identify the full commit and original source status; deployment files also retain source metadata. The training, compilation, evaluation, and analysis entry points are \texttt{train\_witt.py}, \texttt{prepare\_gate\_tables.py}, \texttt{evaluate\_gate\_tables\_paired.py}, and \texttt{analyze\_paired\_gate\_results.py}, respectively. Plot-generation scripts read saved JSON results and require no model inference. The local package contains configuration files, raw metrics, figure sources, environment specifications, and a reproduction guide. Public hosting and author details remain pending for this course draft.

\paragraph{Saved evidence.}
The selected checkpoint is identified by \texttt{R003-checkpoint-selection.json}. P002's preparation record supplies parameter counts and complete file sizes. P002 and P003 retain validation metrics; P004 retains per-method, per-noise, per-batch error sums and counts on the official test images. The test-freeze record precedes testing. Source snapshots accompanying P002--P004 preserve the evaluator and protocol. Original dataset pixels are separate from these source snapshots.

\paragraph{Actual AI use.}
The project follows the ARIS skill-based research workflow~\cite{yin2026aris}, with Codex as the executing AI. The workflow covers baseline selection, author-code inspection, resumable local training, candidate generation, critical review, method refinement, gated evaluation, and evidence review. Initial work executed several subflows; a user correction led to explicit main-pipeline tracking and the later review--revision--evaluation cycle. Reviewer identities and raw replies are recorded. The reviewers belong to the same model family as the executor, so their assessments remain provisional.

One critical review corrected the three-node pilot's full-grid worst loss from a narrower-subset value to 0.093756 dB, identified missing fractional SNR conditions, and motivated the stronger validation protocol. Code review then required a synthetic serialization witness for testing and verification of actual deployment sizes and checkpoint identity before opening test images. Both fixes preceded P004. A final joint evidence review checked experimental integrity and bounded claims; it was one review session serving several checklists, rather than multiple independent reviews. The resulting artifact is a CVPR-formatted course study. Student-authored reflections and publication decisions are recorded separately from measured scientific results.
